%=============================================================================!
% 0.3  HOW TO USE THE BOOK                                                    !
%=============================================================================!
\section{How to use the book}
\label{sec:or-use}

The book assumes school mathematics and nothing else: algebra,
trigonometry, and the differentiation and integration of a function of one
variable. It assumes no physics. Every further tool, from vectors and
matrices to differential equations, partial derivatives and probability,
is explained where it is first needed. Part~I
starts from what a velocity is and builds, one idea at a time, all the
classical mechanics that molecular dynamics uses; each idea is first seen
in an everyday system, a thrown ball or a block on a spring, and then
applied to atoms by the same argument. Statistical mechanics, molecular
simulation and machine learning are built in the same way from their
foundations. The notebook demonstrations and chapter exercises also allow
the material to be used for teaching; longer sections can be spread over
several sessions.

Each chapter has a companion notebook with the same sections. The book
carries the argument and the mathematics; the notebook carries the same
ideas as code to run and vary, and a pointer in the text says what to
explore. No programming is assumed: Notebook~01 opens with the little
Python that the early notebooks use.

Make a working copy of each notebook, run its cells in order, then return
to an example and change a parameter whose effect you can predict. Compare
the result with that prediction and record the reason for any difference.
The exercise cells leave space for your answers, with worked solutions
folded below them. These calculations are intended to be revisited as new
methods give another way to check an earlier result.

\notebook{01}{the thrown ball: move the time cursor, watch $x$, $v$ and $a$}

The code grows chapter by chapter in a small package, \texttt{mdlab}, from
a single thrown ball to a molecular dynamics engine with thermostats
and analysis tools, then to the regression and optimisation routines of
Chapter~18. The planned continuation develops equivariant networks.
Every listing is taken from it, and every physical routine in it is tested
against an analytic result or an independent implementation.

Four boxes set material apart from the argument.

\begin{keyidea}
The result of a section that later chapters rely on.
\end{keyidea}

\begin{bridge}
A connection to electronic-structure theory, for readers who know it. The
argument never depends on these boxes, and a reader new to the subject can
pass over them.
\end{bridge}

\begin{practice}
What goes wrong in real simulations, and how to diagnose it.
\end{practice}

\begin{toolbox}
A mathematical tool, such as the Taylor series, derived in full where it is
first needed. A reader who already knows it can pass over the box.
\end{toolbox}
